\chapter*{Introduction}
\addcontentsline{toc}{chapter}{Introduction}
\markboth{Introduction}{Introduction}

Le stage opérateur de première année n'a pas pour objet de réaliser un travail d'ingénieur, mais de découvrir, de l'intérieur, le fonctionnement d'une organisation et le monde professionnel. C'est dans cet esprit que j'ai effectué mon stage au sein de Reminex, la filiale d'ingénierie du Groupe Managem, rattachée à la Direction Ingénierie.

Je suis arrivée avec une connaissance très générale de l'industrie et une idée encore floue de ce que recouvre concrètement le métier d'ingénieur dans un grand groupe. Mon objectif était donc avant tout d'observer et de comprendre~: comment s'organise une entreprise minière de dimension africaine~? Comment naît et se pilote un projet industriel, depuis l'exploration jusqu'à la production~? Qu'est-ce que l'ingénierie de procédé, dont j'avais surtout entendu parler en cours de manière abstraite~? Et comment les différents métiers d'une même entreprise communiquent-ils entre eux~?

Pour répondre à ces questions, mon stage a reposé sur deux types d'activités complémentaires. D'une part, j'ai contribué, sous la supervision de mon tuteur, à la préparation de rapports hebdomadaires de suivi d'un chantier réel~: l'installation des convoyeurs de la mine de cuivre de Tizert. D'autre part, j'ai bénéficié de plusieurs présentations conduites par des ingénieurs expérimentés, qui m'ont permis de relier ce chantier à une vision plus large de l'ingénierie minière.

Ce rapport rend compte de cette découverte. Il présente d'abord le cadre dans lequel s'est déroulé le stage --- le Groupe Managem et sa filiale Reminex ---, puis la chaîne de valeur minière telle qu'elle m'a été expliquée et les échanges que j'ai eus avec les ingénieurs. Il décrit ensuite les projets que j'ai découverts, le déroulement concret de mon stage et les tâches qui m'ont été confiées. Il se termine par une analyse personnelle de ce que cette expérience m'a apporté et de la manière dont elle a fait évoluer ma vision de l'ingénierie.

Conformément aux consignes de rédaction, les informations factuelles sur le Groupe et ses projets, issues de sources publiques, sont référencées et clairement distinguées de ce que j'ai personnellement observé ou de ce qui m'a été présenté oralement pendant le stage.
